Уявіть, що одного ранку ви прокидаєтеся, і, розплющивши очі, раптом усвідомлюєте, що все вже не так, як було. І я маю на увазі все! Не те, що щось інакше виглядає для ока, а те, що те, що дивиться, якимось дивним і непередбачуваним чином змінилося чи переродилося. І ось ви, щойно ступивши з ліжка, дивуєтеся, чи не спите ви насправді ще під ковдрою й не сните цей дивний сон.

Але ви не спите, і ви це знаєте. Ви знаєте, що не спите, з такою певністю, якої ніколи раніше не відчували. Насправді ви усвідомлюєте, що кожна мить кожного дня вашого життя аж до кількох хвилин тому минула в стані несвідомого сновидіння наяву, яке здавалося настільки реальним і справжнім, що вам навіть не спадало на думку засумніватися в його дійсності. Загальна думка така: щойно ви встали з ліжка й ходите, ви цілком прокинулися. Але це просто неправда.

Як ви не помітили чогось настільки очевидного?

Але це ще не все, і навіть близько не все. Уявіть, що, помічаючи цю дивну зміну в собі, ви усвідомлюєте дещо ще бентежніше: що всередині вас немає нічого; насправді немає жодного «я», щоб бути всередині. Тож ви починаєте шукати себе. Як вам узагалі вдалося загубити себе? Зрештою, ви не пара капців, які можна просто десь забути. Ви — це ви, і ви мусите бути десь тут. Хіба не так?

Але хоч би де ви заглядали всередину, ви не можете знайти себе ні як когось, ні в якомусь певному місці. Усі старі думки й усі старі спогади більше не стосуються вас самих. Вони порожні — порожні щодо «я». Навіть те обличчя в дзеркалі, що здається таким знайомим, без «я». Воно ніколи його не мало, крім того, яке ви вигадали у своєму розумі. І навіть те вигадане ніколи насправді не мало «я».

Ви дивитеся у вікно і... немає жодного «там, зовні». Скрізь, куди ви дивитеся, усе якось усередині — хоч би де було це «всередині». Тож ви дивитеся крізь вікно, і все там, чи он там, — усередині вас, і не просто всередині вас, а й є вами. Земля, і дерева, і паркан он там, а також небо над головою і білі пухнасті хмари — усе це ви. Це не має сенсу, але це очевидно, як дихання. Що це за «я», яке те саме, що й усе інше?

Це справді дивно.

Уявіть, що ви йдете тротуаром і слухаєте, як сусіди розмовляють одне з одним, коли вам спадає на думку, що вони все це вигадують. Усі історії, усі дрібні судження, усі твердо засвоєні погляди, усі «він мав би», «вона мала би» й «я вважаю, що...» — усе вигадане, але сприймається як реальне. Наче вони граються в уявне, але забувають, що це все уява, чи губляться в учорашньому сні.

Як вони можуть ставитися до того, що кажуть, так серйозно, ніби це має значення або має якусь основу в Дійсності? Як вони не бачать? Але вони не бачать. Для них це їхня дійсність — єдина, яку вони знають і, мабуть, знатимуть. Як же це дивно.

Тепер уявіть, що ви зупиняєтеся й сідаєте на лавку в парку. Коли ви там сидите, усе зупиняється, цілком зупиняється. Ваш розум такий тихий і нерухомий, що ви чуєте, як у повітрі плавають порошинки. Раптом ви падаєте, падаєте й падаєте. Немає ні землі внизу, ні неба вгорі — лише громова тиша, що трощить, і дедалі шаленіша швидкість. Ви раптом усвідомлюєте, що це вас уб'є, розірве на шматки й оберне ваші легені на порох. Немає виходу, немає жодного способу вижити. І ви робите єдине, що можна зробити.

Віддання.

Усе стає білим і порожнім, порожнішим за безмежний простір.

Перед життям і смертю ви зникаєте з існування (чи, може, входите в нього?). Позачасовість — це все, що є, усе, що будь-коли було чи могло бути. Вічність панує, і вона променисто присутня в кожній часточці буття.

Щось непороджене й невмируще оживає й розплющує очі — ваші очі. Ви або Воно досі сидить на парковій лавці. Воно усміхається, променисте й задоволене. Повз проходить дівчинка на роликах. Сонце блищить крізь листя осик, а старий чоловік курить люльку на пішохідному містку над струмком, що впадає в ставок, повний золотих рибок.

Скрізь, куди ви дивитеся, — порожнеча. Кожна «річ» — це завіса, покров, що закриває Безкінечне. Ніщо не таке, як здається, і все — точно таке, як є. Якось досконале в усьому своєму уявному хаосі, Безкінечне панує. Ви знаєте з точною певністю, що немає нічого іншого — ніщо не могло б бути іншим, ніж ця величезна й абсолютна порожнеча, цей чистий і безкінечний Потенціал, це непороджене й неоформлене Безкінечне.

Ви озираєте своє життя й усвідомлюєте, що все, що будь-коли сталося чи могло статися, — від народження, крізь усі злети й падіння цього швидкоплинного життя, крізь дивні здійснення духовного пробудження, аж до цієї точної миті поза часом, — було і є миттєвим виявом, справді лише спалахом, безмежного потенціалу Безкінечного, що входить в існування й виходить із нього.

Стара подруга знаходить вас на лавці в парку. Вона сідає поруч і питає: «Чим ти займаєшся?» Ви любите її, як люблять друзів, але що ви можете сказати? Ви вже без слів, і всередині так само тихо, як у мертвих. Вона цього не знає, але ви в двох різних світах, які дивно перетинаються тут, на цій парковій лавці. Як перетнути безкінечність, щоб порозумітися з нею?

Якусь мить ви напружуєтеся всередині, шукаючи слів для відповіді. Настає мовчазна пауза — чи вона щось запідозрила? Чи відчуває, що щось не так? Прохолодний вітерець пестить ваше обличчя, і всесвіт усміхається у вас. «Та ні, нічого», — кажете ви. «Абсолютно нічого».
